\documentclass{article}
\usepackage{amsmath}
\usepackage{amsfonts}
\usepackage{url}
\usepackage{graphicx}

\usepackage{accessibility}
\usepackage{amssymb}

\title{Mathematical Foundations of Clinical Trial Statistics: An Implementation Perspective}
\author{Jules (AI Agent)}
\date{\today}

\begin{document}

\maketitle

\begin{abstract}
Clinical trials are the backbone of evidence-based medicine, relying heavily on rigorous statistical methods to discern true treatment effects from random noise. This paper documents the mathematical principles underlying the \texttt{clinical\_trials} module in the \texttt{math\_explorer} Rust library. We cover hypothesis testing, randomization strategies, sample size estimation, and survival analysis, providing the theoretical basis for the implemented algorithms.
\end{abstract}

\section{Introduction}
The primary goal of a clinical trial is to evaluate the efficacy and safety of a new intervention. To ensure scientific validity, trials must be designed and analyzed using robust statistical frameworks. We have implemented a suite of tools covering the lifecycle of a trial, from design (randomization) to planning (power analysis) and final analysis (efficacy and survival).

\section{Hypothesis Testing}
At the core of clinical inference is hypothesis testing. We test the Null Hypothesis ($H_0$) against an Alternative Hypothesis ($H_1$).

\subsection{T-Test for Continuous Outcomes}
To compare the means of two independent groups (e.g., Treatment vs. Control), we use the independent samples t-test. The t-statistic is calculated as:
\begin{equation}
t = \frac{\bar{X}_1 - \bar{X}_2}{s_p \sqrt{\frac{1}{n_1} + \frac{1}{n_2}}}
\end{equation}
where $s_p^2$ is the pooled variance:
\begin{equation}
s_p^2 = \frac{(n_1 - 1)s_1^2 + (n_2 - 1)s_2^2}{n_1 + n_2 - 2}
\end{equation}
The p-value is derived from the Student's t-distribution with $df = n_1 + n_2 - 2$.

\subsection{Chi-Square Test for Categorical Outcomes}
For binary outcomes (e.g., Cured/Not Cured), we employ the Pearson Chi-Square test on a $2 \times 2$ contingency table:
\begin{equation}
\chi^2 = \sum \frac{(O_i - E_i)^2}{E_i}
\end{equation}
where $O_i$ is the observed frequency and $E_i$ is the expected frequency under independence.

\section{Study Design: Randomization}
To minimize selection bias, patients are assigned to groups using random algorithms.
\begin{itemize}
    \item \textbf{Simple Randomization}: Independent Bernoulli trials ($p=0.5$).
    \item \textbf{Block Randomization}: Ensuring balance within small blocks of size $k$. This prevents severe imbalance in group sizes.
    \item \textbf{Stratified Randomization}: Performing randomization separately within strata (e.g., age groups) to ensure prognostic factors are balanced.
\end{itemize}

\section{Sample Size Calculation}
Adequate sample size is crucial to achieve sufficient statistical power ($1-\beta$).

\subsection{Comparing Means}
The formula for the sample size per group ($n$) to detect a difference $\Delta$ with standard deviation $\sigma$:
\begin{equation}
n = \frac{2\sigma^2 (Z_{1-\alpha/2} + Z_{1-\beta})^2}{\Delta^2}
\end{equation}

\subsection{Comparing Proportions}
For two proportions $p_1$ and $p_2$:
\begin{equation}
n = \frac{(p_1(1-p_1) + p_2(1-p_2)) (Z_{1-\alpha/2} + Z_{1-\beta})^2}{(p_1 - p_2)^2}
\end{equation}

\section{Analysis Metrics}
\subsection{Relative Risk and Odds Ratio}
\begin{align}
RR &= \frac{a / (a+b)}{c / (c+d)} \\
OR &= \frac{ad}{bc}
\end{align}
Confidence intervals are calculated using the standard error of the natural log of these estimates.

\section{Survival Analysis}
For time-to-event data, we implement the Kaplan-Meier estimator:
\begin{equation}
\hat{S}(t) = \prod_{t_i \le t} \left(1 - \frac{d_i}{n_i}\right)
\end{equation}
where $d_i$ is the number of events at time $t_i$ and $n_i$ is the number at risk.

\bibliographystyle{plain}
\bibliography{references}

\end{document}
